% GENERATED FILE. Edit canonical lessons, then run npm run generate:books.
% canonical-source-hash: fnv1a64:218e893553f92da8
% canonical-lessons: ES-C709-aislamiento, ES-C709-solitario, ES-C709-extranar, ES-C709-nostalgia, ES-C709-anorar

\chapter{Echar de menos}
\label{ch:es-echar-de-menos}
\hlchaptermodality{709}

\begin{chapteropening}
\textbf{By the end of this chapter:} \emph{I can say that I feel lonely and miss people and places.}

Complete the last lesson of chapter 709: I can say that I feel lonely and miss people and places.
\end{chapteropening}

\section[el aislamiento]{el aislamiento — isolation}

\label{lesson:ES-C709-aislamiento}



\begin{quote}
Before the new one: say the Spanish for to envy, then the Spanish for resentment.
\end{quote}

\subsection*{You'll want to know: el aislamiento}
\textbf{el aislamiento} — \textquotedblleft{}isolation\textquotedblright{}.

\begin{quote}
Muchas personas mayores viven en aislamiento.
\end{quote}

\begin{quote}
\emph{Many older people live in isolation.}
\end{quote}

\emph{Aislado} — isolated, cut off.

\begin{cousinweb}
\emph{aislamiento} is built on \emph{aislar} \textquotedblleft{}to isolate\textquotedblright{}, made from \emph{isla} \textquotedblleft{}island\textquotedblright{}, which is inherited from Latin \emph{īnsula} (EtymDB): to be cut off like an island.
\end{cousinweb}

\begin{grammarlens}[title={What you've built}]
One more word for telling, explaining and describing at B1.
\end{grammarlens}

\subsection*{Your turn}
\begin{itemize}
\raggedright
  \item \emph{Say it:} \emph{el aislamiento}
  \item \emph{Say it:} \emph{el aislamiento}, once more
  \item \emph{Say it:} a sentence of your own with \emph{el aislamiento}
  \item \emph{Recall:} say the Spanish for to envy, then the Spanish for resentment, then say \emph{el aislamiento} again
\end{itemize}

\begin{tcolorbox}[breakable,colback=teal!4,colframe=teal!35!black,title={Before you move on}]
What does el aislamiento mean? (\textquotedblleft{}Isolation\textquotedblright{}.) Say it once more.
\end{tcolorbox}
\section[solitario]{solitario — lonely, solitary}

\label{lesson:ES-C709-solitario}



\begin{quote}
Before the new one: say the Spanish for resentment, then the Spanish for isolation.
\end{quote}

\subsection*{You'll want to know: solitario}
\textbf{solitario} — \textquotedblleft{}lonely, solitary\textquotedblright{}.

\begin{quote}
Es un hombre solitario.
\end{quote}

\begin{quote}
\emph{He is a solitary man.}
\end{quote}

Also a card game you play alone: \emph{hacer un solitario}.

\begin{cousinweb}
\emph{solitario} is a learned borrowing from Latin \emph{sōlitārius} \textquotedblleft{}lonely\textquotedblright{} (EtymDB); the same Latin word, inherited, gave \emph{soltero}.
\end{cousinweb}

\begin{grammarlens}[title={What you've built}]
One more word for telling, explaining and describing at B1.
\end{grammarlens}

\subsection*{Your turn}
\begin{itemize}
\raggedright
  \item \emph{Say it:} \emph{solitario}
  \item \emph{Say it:} \emph{solitario}, once more
  \item \emph{Say it:} a sentence of your own with \emph{solitario}
  \item \emph{Recall:} say the Spanish for resentment, then the Spanish for isolation, then say \emph{solitario} again
\end{itemize}

\begin{tcolorbox}[breakable,colback=teal!4,colframe=teal!35!black,title={Before you move on}]
What does solitario mean? (\textquotedblleft{}Lonely, solitary\textquotedblright{}.) Say it once more.
\end{tcolorbox}
\section[extrañar]{extrañar — to miss (someone); to find strange}

\label{lesson:ES-C709-extranar}



\begin{quote}
Before the new one: say the Spanish for isolation, then the Spanish for lonely.
\end{quote}

\subsection*{You'll want to know: extrañar}
\textbf{extrañar} — \textquotedblleft{}to miss (someone); to find strange\textquotedblright{}.

\begin{quote}
Te extraño mucho.
\end{quote}

\begin{quote}
\emph{I miss you a lot.}
\end{quote}

To miss is the Latin American sense; in Spain people mostly say \emph{echar de menos}. \emph{Me extraña} — I find it strange.

\begin{cousinweb}
\emph{extrañar} is inherited from Late Latin \emph{extrāneāre}, built on \emph{extrāneus} \textquotedblleft{}external, foreign, strange\textquotedblright{} (EtymDB); English \emph{strange} comes from the same adjective through French.
\end{cousinweb}

\begin{grammarlens}[title={What you've built}]
One more word for telling, explaining and describing at B1.
\end{grammarlens}

\subsection*{Your turn}
\begin{itemize}
\raggedright
  \item \emph{Say it:} \emph{extrañar}
  \item \emph{Say it:} \emph{extrañar}, once more
  \item \emph{Say it:} a sentence of your own with \emph{extrañar}
  \item \emph{Recall:} say the Spanish for isolation, then the Spanish for lonely, then say \emph{extrañar} again
\end{itemize}

\begin{tcolorbox}[breakable,colback=teal!4,colframe=teal!35!black,title={Before you move on}]
What does extrañar mean? (\textquotedblleft{}To miss (someone); to find strange\textquotedblright{}.) Say it once more.
\end{tcolorbox}
\section[la nostalgia]{la nostalgia — nostalgia, homesickness}

\label{lesson:ES-C709-nostalgia}



\begin{quote}
Before the new one: say the Spanish for lonely, then the Spanish for to miss.
\end{quote}

\subsection*{You'll want to know: la nostalgia}
\textbf{la nostalgia} — \textquotedblleft{}nostalgia, homesickness\textquotedblright{}.

\begin{quote}
Siento nostalgia de mi pueblo.
\end{quote}

\begin{quote}
\emph{I feel homesick for my village.}
\end{quote}

\emph{Nostalgia de} + place or time.

\begin{cousinweb}
\emph{nostalgia} was borrowed from New Latin \emph{nostalgia} (EtymDB), a word a Swiss doctor coined from Greek \emph{nóstos} \textquotedblleft{}return home\textquotedblright{} and \emph{álgos} \textquotedblleft{}pain\textquotedblright{}.
\end{cousinweb}

\begin{grammarlens}[title={What you've built}]
One more word for telling, explaining and describing at B1.
\end{grammarlens}

\subsection*{Your turn}
\begin{itemize}
\raggedright
  \item \emph{Say it:} \emph{la nostalgia}
  \item \emph{Say it:} \emph{la nostalgia}, once more
  \item \emph{Say it:} a sentence of your own with \emph{la nostalgia}
  \item \emph{Recall:} say the Spanish for lonely, then the Spanish for to miss, then say \emph{la nostalgia} again
\end{itemize}

\begin{tcolorbox}[breakable,colback=teal!4,colframe=teal!35!black,title={Before you move on}]
What does la nostalgia mean? (\textquotedblleft{}Nostalgia, homesickness\textquotedblright{}.) Say it once more.
\end{tcolorbox}
\section[añorar]{añorar — to long for, to yearn for}

\label{lesson:ES-C709-anorar}



\begin{quote}
Before the new one: say the Spanish for to miss, then the Spanish for nostalgia.
\end{quote}

\subsection*{You'll want to know: añorar}
\textbf{añorar} — \textquotedblleft{}to long for, to yearn for\textquotedblright{}.

\begin{quote}
Añoro los veranos de mi infancia.
\end{quote}

\begin{quote}
\emph{I long for the summers of my childhood.}
\end{quote}

A little literary; it takes a direct object: \emph{añoro mi casa}.

\begin{cousinweb}
\emph{añorar} was borrowed from Catalan \emph{enyorar}, which comes from Latin \emph{ignōrāre} \textquotedblleft{}not to know\textquotedblright{} (EtymDB); the sense moved from not knowing where someone is to missing them.
\end{cousinweb}

\begin{grammarlens}[title={What you've built}]
One more word for telling, explaining and describing at B1.
\end{grammarlens}

\subsection*{Your turn}
\begin{itemize}
\raggedright
  \item \emph{Say it:} \emph{añorar}
  \item \emph{Say it:} \emph{añorar}, once more
  \item \emph{Say it:} a sentence of your own with \emph{añorar}
  \item \emph{Recall:} say the Spanish for to miss, then the Spanish for nostalgia, then say \emph{añorar} again
\end{itemize}

\begin{tcolorbox}[breakable,colback=teal!4,colframe=teal!35!black,title={Before you move on}]
What does añorar mean? (\textquotedblleft{}To long for, to yearn for\textquotedblright{}.) Say it once more.
\end{tcolorbox}
